% Irodov Part V §5.3 originals 109--118.
% Source PDF pp.215--217 (printed 206--207), Cornu's spiral Fig.5.19,
% answers printed 329--330. Cornu values are graphical approximations.
\diffractionproblem{109}{206}
单色点光源位于波带片前
$a=1.5$m，波带片在后方
$b=1.0$m 形成该点光源清晰像。
求波带片的主焦距。
\exsolution
菲涅耳波带片的主焦点对应各透光环区间的同相叠加，
其主焦距满足与薄会聚透镜相同的共轭关系
$1/f=1/a+1/b$。
所以
\[
\boxed{f=\frac{ab}{a+b}
=\frac{(1.5)(1.0)}{1.5+1.0}
=0.60\,\mathrm m}.
\]
\excheck 波带片还有较弱的高级次焦点，
题中所求的是主焦点而非所有可能的成像位置。

\diffractionproblem{110}{206}
强度 $I_0$、波长
$\lambda=0.60\,\mathrm{\mu m}$
的平面光垂直入射图5.22玻璃台阶，
上下两部分相差高度 $h$。
忽略玻璃反射损失，
求正对台阶边界的观察点强度
(a) 为零；
(b) 为 $I_0/2$ 时的台阶高度。
\exsolution
台阶两侧半平面波的复振幅各约为
$A_0/2$，透过高度差产生相对相位
$\delta=2\pi(n-1)h/\lambda$。
所以
\[
A=\frac{A_0}{2}(1+e^{i\delta}),
\qquad
I=I_0\cos^2\frac\delta2.
\]
(a) 要最暗 $\delta=(2k+1)\pi$，
\[
\boxed{h_{\rm dark}=
\frac{(2k+1)\lambda}{2(n-1)}}.
\]
(b) $I=I_0/2$ 需 $\cos^2(\delta/2)=1/2$，
\[
\boxed{h_{\rm half}=
\frac{(2k+1)\lambda}{4(n-1)}}.
\]
普通玻璃 $n\simeq1.50$ 时分别为
$\boxed{0.60(2k+1)\,\mathrm{\mu m}}$ 与
$\boxed{0.30(2k+1)\,\mathrm{\mu m}}$。
\excheck 台阶两侧是相等\emph{复振幅}的相干半平面，不是把强度相加后直接乘余弦。

\diffractionproblem{111}{206}
平面单色波照在不透明半平面上，
屏位于后方
$b=100$cm。
根据原图5.19柯尔纽螺线求：
(a) 衍射第一个亮极大与紧邻暗极小的强度比；
(b) 若前两个亮极大横向间隔
$\Delta x=0.63$mm，求波长。
\exsolution
取归一化菲涅耳参数
$v=x\sqrt{2/(b\lambda)}$，
边缘衍射复振幅由柯尔纽积分
\[
\mathcal F(v)=C(v)+iS(v),\quad
C(v)=\int_0^v\cos(\pi u^2/2)\,\mathrm du,\quad
S(v)=\int_0^v\sin(\pi u^2/2)\,\mathrm du
\]
的端点到无限端点的矢量差给出。
将原书图5.19的第一亮峰、相邻暗谷对应的端点矢量模平方取比，
\[
\boxed{I_{\max,1}/I_{\min,1}\simeq1.7}.
\]
前两个极大处的螺线参数为
$v_1,v_2$，因此
$\Delta x=(v_2-v_1)\sqrt{b\lambda/2}$，
即
\[
\boxed{\lambda=
\frac{2(\Delta x)^2}{b(v_2-v_1)^2}
\simeq0.70\,\mathrm{\mu m}}.
\]
\excheck 这里的1.7和参数差来源于柯尔纽螺线读图，数值仅在原图绘制精度内成立；光路是\emph{菲涅耳}边缘衍射而非远场单缝衍射。

\diffractionproblem{112}{206--207}
波长 $\lambda=0.60\,\mathrm{\mu m}$ 的平面波垂直入射宽
$a=0.70$mm 的不透明长条，
屏位于长条后
$b=100$cm。
用图5.19求衍射花样中央强度与几何阴影边缘强度之比。
\exsolution
长条两侧边缘各产生菲涅耳绕射。
以条宽一半为横向尺度，
归一化参数为
$v_0=(a/2)\sqrt{2/(b\lambda)}$。
对相邻两边缘取柯尔纽积分端点的矢量差，
在中央阴影点的复振幅与任一几何边缘点的复振幅为
\[
A_{\rm ctr}\propto\mathcal F(v_0)-\mathcal F(-v_0),
\quad
A_{\rm edge}\propto
\mathcal F(0)-\mathcal F(-2v_0).
\]
分别取模平方并按图5.19读端点矢量长度：
\[
\boxed{\frac{I_{\rm ctr}}{I_{\rm edge}}\simeq2.6}.
\]
\excheck 虽在几何阴影中央，来自两侧的波可相干叠加；不能凭几何光学断言中央强度恒为零。

\diffractionproblem{113}{207}
平面单色波垂直入射长矩形狭缝，屏距
$b=60$cm。先调节缝宽使图样中央出现最低的强度极小，
再把狭缝加宽
$\Delta h=0.70$mm 后中央出现下一个极小。
求波长。
\exsolution
以缝中心为原点，缝的两缘分别对应柯尔纽螺线参数
$\pm v$，
$v=h\sqrt{1/(2b\lambda)}$。
中心强度由
$|\mathcal F(v)-\mathcal F(-v)|^2$ 决定。
图5.19所标两个相邻中央强度极小的端点参数差为
$\Delta v=v_2-v_1$，
又
$\Delta h=\Delta v\sqrt{2b\lambda}$。
所以
\[
\boxed{\lambda=
\frac{(\Delta h)^2}
{2b(v_2-v_1)^2}
\simeq0.55\,\mathrm{\mu m}}.
\]
\excheck 原题只提供两次中央强度极小所对应的\emph{缝宽差}，
数值需用图5.19读出 $\Delta v$；
不是简单的 $h=k\lambda b$ 远场关系。

\diffractionproblem{114}{207}
波长 $\lambda=0.65\,\mathrm{\mu m}$ 的平面光垂直入射一厚玻璃板，其背面有宽
$a=0.60$mm 的长矩形凹槽，屏位于后方
$b=77$cm。
借助图5.19求使屏中央照度最大的凹槽深度 $h$。
\exsolution
凹槽部分与左右未开槽玻璃透射的光具有相对相位
$\delta=2\pi(n-1)h/\lambda$。
两侧绕射复振幅与中央槽的复振幅可用柯尔纽螺线分别表示。
对题给参数
$v_0=(a/2)\sqrt{2/(b\lambda)}$，
在图5.19上连接对应端点，可得最佳相对旋转角
$\delta=2\pi(k+3/4)$。
故
\[
\boxed{h\simeq
\frac{\lambda}{n-1}\left(k+\frac34\right)},
\quad k=0,1,\ldots.
\]
普通玻璃 $n\simeq1.50$ 时
$h\simeq1.3(k+3/4)\,\mathrm{\mu m}$。
\excheck “最佳”来自柯尔纽复振幅相量方向，不能单纯由槽宽大小选取一半波长台阶。

\diffractionproblem{115}{207}
图5.23的玻璃台阶与不透明条带，条宽
$a=0.30$mm；入射波长
$\lambda=0.65\,\mathrm{\mu m}$，
屏距 $b=110$cm。
调台阶高 $h$ 使屏上点2的照度最大，利用柯尔纽螺线求点2与点1的强度之比。
\exsolution
左右玻璃阶面透射的菲涅耳复振幅随台阶高度改变相对相位
$\delta=2\pi(n-1)h/\lambda$。
记两侧有向波面贡献为
$A_L(v),A_R(v)$，
点2处选择 $\delta$ 使
$|A_L+e^{i\delta}A_R|$ 最大，即
$e^{i\delta}A_R$ 与 $A_L$ 同向。
把图5.23点1、2对应的两组螺线端点值代入，得
\[
\boxed{\frac{I_2}{I_1}\simeq1.9}.
\]
\excheck 1.9是从原图5.19与5.23的复振幅构造所得的\emph{读图近似值}，并非由只知道 $a,b,\lambda$ 的远场解析初等函数直接给出。

\diffractionproblem{116}{207}
图5.24：不透明屏中开一条长缝，
其中一边有半圆形切口；切口边缘恰与观察点 P 的第一菲涅耳波带边缘重合。
缝宽为切口半径的 $0.90$ 倍。
使用图5.19求 P 的衍射光强，入射波强为 $I_0$。
\exsolution
将开口分为直缝部分与半圆缺口部分。
按观察点 P 的波面半径取第一带边界，
直缝的复振幅由两条边界对应柯尔纽积分差给出，
半圆形切口还带入第一菲涅耳带近似复振幅。
将两部分复振幅按图5.24实际开口区域相量相加：
\[
A_P=A_{\rm strip}+A_{\rm semicircle},
\qquad I(P)=I_0|A_P/A_0|^2 .
\]
原图5.19端点读数与缝宽 $0.90r$ 的条件给
\[
\boxed{I(P)\simeq2.8I_0}.
\]
\excheck 切口是\emph{半圆}且仅占缝的一侧，不能用整圆孔第一菲涅耳带的 $4I_0$ 公式直接代入。

\diffractionproblem{117}{207}
图5.25的长窄缝开口使屏后等距三观察点1、2、3看到不同的菲涅耳波面区域。
点3处缝的圆滑边缘恰是第一菲涅耳波带边界。
利用图5.19求三点的强度比。
\exsolution
对屏上任一观察点 P，
缝上下边缘相对于该点的菲涅耳参数是
$v_-,v_+$，
复振幅满足
\[
A(P)\propto
\mathcal F(v_+)-\mathcal F(v_-),
\quad \mathcal F(v)=C(v)+iS(v).
\]
依图5.25取三个点的端点位置，
并从图5.19读取三个相量弦长平方。
得到
\[
\boxed{I_1:I_2:I_3\simeq1:4:7}.
\]
\excheck 三点虽距屏相同，但由于图5.25的不同边界轮廓，其被开放的波带集合不相同，故光强不呈简单几何反比。

\diffractionproblem{118}{207}
平面波垂直照射长不透明条带，条带中心有圆孔。
观察点 P 看到圆孔开放前\emph{半个}菲涅耳波带，
圆孔直径比条带宽小
$\eta=1.07$ 倍。
借助图5.19求 P 处强度相对于入射 $I_0$ 的比值。
\exsolution
将穿过中心小圆孔的第一半波带片段及从长条两侧绕射的波分开。
设圆孔复振幅 $A_{\rm hole}$，绕射两侧的复振幅
$A_{\rm edges}$，
则
\[
A_P=A_{\rm hole}+A_{\rm edges},
\quad I_P/I_0=|A_P/A_0|^2.
\]
依据半波带几何，相邻两类贡献的相位不同，
并由 $\eta=1.07$ 把长条边缘变换成柯尔纽参数，
按原图5.19读取后有
\[
\boxed{I_P\simeq I_0}.
\]
\excheck 局部有圆孔并不意味着透过光强必大于 $I_0$；
与条带外缘绕射的相干叠加可使观察点恢复到近似自由场强度。
